% Slug: tiny-fault-diagnosis
% Source: chapters/ch01/fault-lab.tex:237 (tikzpicture env start line)
% Original label: \label{fig:tiny-fault-diagnosis}
% Original caption (LaTeX source):
%   The verifier accepts both programs; only one preserves the source value.
% Original caption (plain text):
%   The verifier accepts both programs; only one preserves the source value.
% Required TikZ setup: \usepackage{tikz} + \usetikzlibrary{arrows.meta, positioning, shapes.geometric, fit} (loaded by support/compiler-pseudocode.sty).
% Shared styles used in this picture: `basic block', `cfg edge'.
%   Defined in support/compiler-pseudocode.sty: `basic block', `cfg edge', `back edge', `block label'.
% Usage: % Slug: tiny-fault-diagnosis
% Source: chapters/ch01/fault-lab.tex:237 (tikzpicture env start line)
% Original label: \label{fig:tiny-fault-diagnosis}
% Original caption (LaTeX source):
%   The verifier accepts both programs; only one preserves the source value.
% Original caption (plain text):
%   The verifier accepts both programs; only one preserves the source value.
% Required TikZ setup: \usepackage{tikz} + \usetikzlibrary{arrows.meta, positioning, shapes.geometric, fit} (loaded by support/compiler-pseudocode.sty).
% Shared styles used in this picture: `basic block', `cfg edge'.
%   Defined in support/compiler-pseudocode.sty: `basic block', `cfg edge', `back edge', `block label'.
% Usage: % Slug: tiny-fault-diagnosis
% Source: chapters/ch01/fault-lab.tex:237 (tikzpicture env start line)
% Original label: \label{fig:tiny-fault-diagnosis}
% Original caption (LaTeX source):
%   The verifier accepts both programs; only one preserves the source value.
% Original caption (plain text):
%   The verifier accepts both programs; only one preserves the source value.
% Required TikZ setup: \usepackage{tikz} + \usetikzlibrary{arrows.meta, positioning, shapes.geometric, fit} (loaded by support/compiler-pseudocode.sty).
% Shared styles used in this picture: `basic block', `cfg edge'.
%   Defined in support/compiler-pseudocode.sty: `basic block', `cfg edge', `back edge', `block label'.
% Usage: % Slug: tiny-fault-diagnosis
% Source: chapters/ch01/fault-lab.tex:237 (tikzpicture env start line)
% Original label: \label{fig:tiny-fault-diagnosis}
% Original caption (LaTeX source):
%   The verifier accepts both programs; only one preserves the source value.
% Original caption (plain text):
%   The verifier accepts both programs; only one preserves the source value.
% Required TikZ setup: \usepackage{tikz} + \usetikzlibrary{arrows.meta, positioning, shapes.geometric, fit} (loaded by support/compiler-pseudocode.sty).
% Shared styles used in this picture: `basic block', `cfg edge'.
%   Defined in support/compiler-pseudocode.sty: `basic block', `cfg edge', `back edge', `block label'.
% Usage: \input{graphs-code/tiny-fault-diagnosis.tikz.tex} inside a figure environment
%   (wrap with \begin{figure} ... \caption{...} \label{fig:tiny-fault-diagnosis} ... \end{figure}).
% Note: tikzpicture body below is verbatim from the source; do not reformat.
\begin{tikzpicture}[node distance=13mm and 17mm]
    \node[basic block] (source) {$0\times1$};
    \node[basic block,below left=of source] (good) {PUSH 0; PUSH 1\\MUL8};
    \node[basic block,below right=of source] (bad) {PUSH 0; PUSH 1\\ADD8};
    \node[basic block,below=of good] (zero) {valid; result 0};
    \node[basic block,below=of bad] (one) {valid; result 1};
    \draw[cfg edge] (source) -- node[above left,font=\small] {correct} (good);
    \draw[cfg edge] (source) -- node[above right,font=\small] {faulty} (bad);
    \draw[cfg edge] (good) -- (zero);
    \draw[cfg edge] (bad) -- (one);
  \end{tikzpicture}
 inside a figure environment
%   (wrap with \begin{figure} ... \caption{...} \label{fig:tiny-fault-diagnosis} ... \end{figure}).
% Note: tikzpicture body below is verbatim from the source; do not reformat.
\begin{tikzpicture}[node distance=13mm and 17mm]
    \node[basic block] (source) {$0\times1$};
    \node[basic block,below left=of source] (good) {PUSH 0; PUSH 1\\MUL8};
    \node[basic block,below right=of source] (bad) {PUSH 0; PUSH 1\\ADD8};
    \node[basic block,below=of good] (zero) {valid; result 0};
    \node[basic block,below=of bad] (one) {valid; result 1};
    \draw[cfg edge] (source) -- node[above left,font=\small] {correct} (good);
    \draw[cfg edge] (source) -- node[above right,font=\small] {faulty} (bad);
    \draw[cfg edge] (good) -- (zero);
    \draw[cfg edge] (bad) -- (one);
  \end{tikzpicture}
 inside a figure environment
%   (wrap with \begin{figure} ... \caption{...} \label{fig:tiny-fault-diagnosis} ... \end{figure}).
% Note: tikzpicture body below is verbatim from the source; do not reformat.
\begin{tikzpicture}[node distance=13mm and 17mm]
    \node[basic block] (source) {$0\times1$};
    \node[basic block,below left=of source] (good) {PUSH 0; PUSH 1\\MUL8};
    \node[basic block,below right=of source] (bad) {PUSH 0; PUSH 1\\ADD8};
    \node[basic block,below=of good] (zero) {valid; result 0};
    \node[basic block,below=of bad] (one) {valid; result 1};
    \draw[cfg edge] (source) -- node[above left,font=\small] {correct} (good);
    \draw[cfg edge] (source) -- node[above right,font=\small] {faulty} (bad);
    \draw[cfg edge] (good) -- (zero);
    \draw[cfg edge] (bad) -- (one);
  \end{tikzpicture}
 inside a figure environment
%   (wrap with \begin{figure} ... \caption{...} \label{fig:tiny-fault-diagnosis} ... \end{figure}).
% Note: tikzpicture body below is verbatim from the source; do not reformat.
\begin{tikzpicture}[node distance=13mm and 17mm]
    \node[basic block] (source) {$0\times1$};
    \node[basic block,below left=of source] (good) {PUSH 0; PUSH 1\\MUL8};
    \node[basic block,below right=of source] (bad) {PUSH 0; PUSH 1\\ADD8};
    \node[basic block,below=of good] (zero) {valid; result 0};
    \node[basic block,below=of bad] (one) {valid; result 1};
    \draw[cfg edge] (source) -- node[above left,font=\small] {correct} (good);
    \draw[cfg edge] (source) -- node[above right,font=\small] {faulty} (bad);
    \draw[cfg edge] (good) -- (zero);
    \draw[cfg edge] (bad) -- (one);
  \end{tikzpicture}
